\section{[INTERNAL] Publication Strategy: ISWC Sequels}
\label{sec:appendix-internal-publication-strategy}
%% ====================================================================
%% INTERNAL appendix --- visible only when compiling draft.tex.
%% Captures the planning discussion (Daniel + Claude, 2026-04-26)
%% on what to do with this work after ISWC.
%% Do NOT include in submission.
%% ====================================================================

\begin{center}
\fbox{\parbox{0.92\textwidth}{\centering\small
\textbf{INTERNAL --- not part of the submission.}\\
This appendix records the publication-strategy discussion behind the
two planned sequels (SWJ + npj Computational Materials) and is
included in the \texttt{draft.tex} build only.}}
\end{center}

\bigskip

The ISWC 2026 paper presents the \MLIPsOntology{} as a Semantic Web
resource. The same infrastructure has natural value for two distinct
audiences, which suggests two journal sequels rather than one.

\subsection{Why Two Sequels?}

The risk of two sequels is salami-slicing: same paper twice, different
framing. The opportunity is genuine dual impact, because the audiences
care about different things:

\begin{itemize}
\item \textbf{Semantic-web / ontology engineers} value transferable
  design patterns: polymorphic-schema generalisation,
  \texttt{rdfs:subPropertyOf} as a non-breaking extension mechanism,
  alignment with established upper ontologies, OWL profile choices,
  SHACL-based validation, SPARQL query patterns.
\item \textbf{Materials scientists who train and use MLIPs} value
  practical metadata standards, FAIR data practices for computational
  materials, integration with existing materials databases (NOMAD,
  Materials Cloud, OPTIMADE), and community recommendations for what
  papers should report. They typically do not read OWL axioms.
\end{itemize}

The two papers can therefore have genuinely different \emph{content
emphases}, not just reframed versions of the same content.

\subsection{Differentiation Plan}

\paragraph{SWJ extension (\texttt{2027-swj-mlips-onto}, this is the
direct extension of the ISWC paper).}
Methodology paper for ontology engineers:
\begin{itemize}
\item Deep dive on the polymorphic-schema and sub-property
  generalisation patterns (already drafted in the SWJ scaffold).
\item Full alignment table with MDO, CMSO, EMMO, ML-Schema, PROV-O,
  Schema.org, QUDT.
\item SHACL-shape catalogue for validation.
\item SPARQL query patterns for cross-module reasoning.
\item Knowledge-graph construction methodology and the Concon platform
  integration.
\item Reuse lessons for other domain ontologies (e.g., neuroscience,
  cheminformatics) that face the same hardcoded-method problem.
\end{itemize}

\paragraph{npj Computational Materials sequel
(\texttt{2027-npj-mlips-onto}, new folder).}
Perspective and data-resource paper for the materials community:
\begin{itemize}
\item Survey of metadata gaps across 10--15 recently published MLIP
  papers; quantitative analysis of what is and is not reported.
\item Multiple encoded case studies (Kumar~2025, Qi~2023, Nitol~2025,
  Gubaev~2023, plus a CCSD(T) example from Yuji's work) showing the
  ontology applied to real published work.
\item FAIR principles applied to MLIP studies; concrete reporting
  recommendations.
\item Integration with NOMAD, Materials Cloud, OPTIMADE.
\item Comparison with materials-data standards (CIF, ChemAxiom).
\item Less heavy on OWL formalism (the MS audience does not read
  axioms); more on practical reporting templates.
\end{itemize}

\subsection{Journal Recommendation for the MS Sequel}

Three candidates considered, in priority order:

\begin{enumerate}
\item \textbf{npj Computational Materials} (Nature). Top-tier in the
  MLIP space, covers data/FAIR/methods, IMW Stuttgart already
  publishes there (Kumar~2025 in Acta Materialia is sister-tier).
  Highest visibility for the MLIP user community. \emph{Primary
  recommendation.}
\item \textbf{Scientific Data} (Nature). Explicitly designed for
  data-resource papers with a FAIR mandate. The \MLIPsOntology{} +
  DaRUS deposit + SHACL shapes is a textbook fit for their format.
  Lower citation prestige than npj Comp Mat but stronger genre
  alignment.
\item \textbf{Digital Discovery} (RSC). Newer journal at the
  chemistry/materials/data intersection, growing rapidly. Less
  established but well-aimed audience.
\end{enumerate}

The recommendation favours npj Computational Materials because (a) the
MLIP community reads it heavily, (b) Blazej's group has standing
there, and (c) reaching the \emph{user} community determines whether
the ontology gets adopted in practice.

%% ====================================================================
